\subsection{The complete Gorenstein curves}\label{forms}
\begin{corollary}\label{list}
Let $R$ be a complete one-dimensional Gorenstein local ring of residue characteristic $p>0$. Then $R$ is weakly normal precisely when one of the following holds.
\begin{enumerate}
\item $R$ is a complete DVR.
\item $R=V_1\times_k V_2$ with complete DVRs of residue field $k$.
\item $R$ is the inverse image of $k$ under $V\to k'$ for a complete DVR $V$ and a separable quadratic field extension $k'/k$.
\end{enumerate}
Under the hypotheses of the converse in Theorem A, these are exactly
the complete one-dimensional Gorenstein local rings satisfying $\K$.
\end{corollary}
\begin{proof}
Weak normality implies seminormality.  Apply
Lemma~\ref{conductor-description}, which requires no separability
assumption, and retain its notation $V,J,k_i$.  Set
\[
 d=\sum_i[k_i:k],\qquad \xi=(\pi_i)_i.
\]
The element $\xi$ is regular, $\xi J=J^2$, and
$\dim_kJ/J^2=d$.  Moreover, $\xi R/J^2$ is the line spanned by the
class of $\xi$.  Hence the maximal ideal of $R/\xi R$ has square
zero and dimension $d-1$.  If $d>1$, this maximal ideal is its socle,
and the Gorenstein property forces $d-1=1$.

For $d=1$, the inverse-image description gives the whole DVR.  For
$d=2$, there are either two branches with residue field $k$, or one
branch with a quadratic residue extension $k'/k$.  If this quadratic
extension were inseparable, its characteristic would be two and
there would be $\alpha\in k'\setminus k$ with $\alpha^2\in k$.
A lift $q\in V$ of $\alpha$ would satisfy $q\notin R$ and $q^2\in R$,
contradicting Lemma~\ref{wn}.  This gives the listed possibilities.

Conversely, in each displayed inverse image the ambient product of DVRs is finite over $R$ because lifts of a residue basis containing $1$ generate it over $R$, as $J$ already belongs to $R$. Eakin's theorem gives Noetherianity, and the componentwise powers of $J$ give completeness. The regular parameter $\xi$ and the socle computation above give the Gorenstein property in the two singular cases. Since $\xi V\subset R$, the total quotient ring is the product of the branch fraction fields. If $q^p\in R$, valuation on each branch puts $q$ in $V$. On the two rational branches equality of residue $p$th powers implies equality of residues. On a separable quadratic branch, an element whose $p$th power lies in $k$ already belongs to $k$. These observations prove root closure, including the DVR case.
\end{proof}

\begin{example}[Two mixed-characteristic nodes]
For every prime $p$, evaluation at $0$ and $p$ identifies
$\mathbf Z_p[z]/(z(z-p))$ with $\mathbf Z_p\times_{\mathbf F_p}\mathbf Z_p$ by $a+bz\mapsto(a,a+pb)$. This ring satisfies $\K$ although its reduction modulo $p$ is $\mathbf F_p[z]/(z^2)$, on which Frobenius has a kernel. The ring $\mathbf Z_p[[t]]/(pt)$ also satisfies $\K$, since evaluation at $t=0$ and reduction modulo $p$ identify it with $\mathbf Z_p\times_{\mathbf F_p}\mathbf F_p[[t]]$. The second example has a vertical branch on which $p$ vanishes.
\end{example}
\begin{example}[The boundary of the quadratic construction]
For $N\ge2$, evaluating at $0,p^N$ identifies $R_N=\mathbf Z_p[z]/(z(z-p^N))$ with the pairs congruent modulo $p^N$. The fraction $q=z/p=(0,p^{N-1})$ is missing, while $pq,q^p\in R_N$ because $p(N-1)\ge N$. For $p=2$ Theorem~\ref{realize} gives a nonsplit module-finite extension. For odd $p$ the determinant bound still restricts the rank of a possible counterexample, but the quadratic construction does not settle $\K$ for $R_N$. Likewise, necessity in mixed characteristic two when two is a nonzero zerodivisor is outside the scope of Theorem A.
\end{example}
